\clearpage
\section{Input safety of micro-beta and reference creation}
\label{sec:copying-operational-safety}
The two operational pre-templates now have full input safety, both in
unrestricted contexts and in the original binding-free context class.
The proof follows the entered-hole consumption argument of
\href{https://arxiv.org/html/1907.01257v5\#A8.SS3.SSS2}{Appendix H.3.2}:
execute the selected left-hand operation, absorb its right-hand graph into
the common context, and retain every other paired filling. The resulting
token is settled. The residual right state represents the original right
state and therefore requires zero machine transitions. This is the
settled-specimen case of input safety, not an assumed
equivalence of the two operational terms.

For micro-beta the exact prefix consists of five search transitions and one
compute transition; for reference creation it consists of three search
transitions and one compute transition. These six/four-step counts concern
the literal operational pre-template graphs. They need not equal the total
cost of a translated term with additional distributors and its continuation.
The original stable-argument premises, lambda-body focus freedom, compute
rules and term translation are retained. Forward and converse robustness
of these two template families remain separate obligations.

\subsection{Exact entry prefixes at the original input}
\termdef{operational-entry-prefix}
An \emph{operational entry prefix} starts with a search token at the unique
term input, performs the original search rules, executes the local operation,
and resumes search at that same input in the right-hand graph.
For reference creation, the first search transition reaches the stable
argument, the second recognises its stable head and returns in backtrack
mode, and the third switches to rewrite mode at the reference operation.
Existence of the stable head and its input incidence follow from the
proved stability lemmas. This argument works for every stable graph,
not only for a constant translation.

\lean{Copying/OperationalInputSafety}{14}{36}{Allocation.entry_trace}

The function \texttt{Kind.searchCost} records five for micro-beta and three
for reference creation. The literal entry-trace theorem uses the existing
micro-beta search trace or the new allocation trace according to the
template kind. It produces an index in the actual input interface, together
with exactly the stated number of search steps.

\lean{Copying/OperationalInputSafety}{41}{43}{Kind.searchCost}
\lean{Copying/OperationalInputSafety}{46}{52}{Literal.entry_trace}

An operational pair may use independent carrier names on its two sides.
Transport the search trace through the inverse source isomorphism and the
local compute transition through both endpoint isomorphisms. Recursive
isomorphisms fix each boundary position, so the starting and resumed input
indices agree. The original literal pair has input interface consisting of
one term wire; the wrapper preserves that interface. Consequently every
entered source position is this same unique input.

\lean{Copying/OperationalInputSafety}{55}{72}{Pair.entry_prefix}
\lean{Copying/OperationalInputSafety}{75}{77}{Pair.input}

\Needspace{11\baselineskip}
\subsection{Absorbing a consumed hole into the common context}
\termdef{consumed-hole-residual}
Let $C$ be the original specimen context, let $e$ be its entered hole,
and let $L_f,R_f$ be its two actual filling families. Write $H=R_e$.
Define the replacement family $K$ by $K_e=H$ and $K_f=L_f$ for $f\ne e$.
The dependent definition transports only the selected edge's interface;
the selected and nonselected equations are proved by the corresponding
equality test. Thus $C[K]$ is precisely the result of replacing this one
component, while the original right-hand state remains $C[R]$.

\lean{Copying/OperationalInputSafety}{83}{89}{OperationalSafety.replace}
\lean{Copying/OperationalInputSafety}{91}{92}{replace_selected}
\lean{Copying/OperationalInputSafety}{94}{95}{replace_other}

The new common context is $C'=C[e\leftarrow H]$, obtained by the actual
single-edge plugging operation. Its remaining filling families assign the
old components to every surviving original edge. Each newly inserted edge
of $H$ receives its own original labelled corolla on both sides. These
families are called \texttt{remaining}; no new specimen constructor is
introduced. The expansion-peeling isomorphism identifies the left result
$C'[\mathrm{remaining}(L)]$ with $C[K]$, and the right result with $C[R]$.
Its host-vertex equation also identifies the resumed token, rather than
merely comparing unpointed graphs.

\lean{Copying/OperationalInputSafety}{99}{108}{OperationalSafety.remaining}
\lean{Copying/OperationalInputSafety}{111}{133}{remaining_iso}

The resumed search token cannot be entering another hole in $C'$. For a
surviving original edge, injectivity of the host inclusion and uniqueness
of an edge sourcing a vertex would force that edge to be $e$, which was
removed. For an inserted edge, a hole label contradicts hole freedom of
$H$. This proves settlement at the actual original input coordinate.
Every surviving edge keeps its old recursive filling certificate: retained
edges, genuinely filled holes and recursively paired boxes are handled
separately. Newly inserted edges have the reflexive corolla certificate.

\lean{Copying/OperationalInputSafety}{137}{155}{consumed_not_entering}
\lean{Copying/OperationalInputSafety}{158}{174}{remaining_edges}

\subsection{The pointed residual specimen and its machine runs}
A finite search trace becomes a run of the combined machine by replacing
each search step with its original combined-search constructor. Induction
preserves the length exactly. For the selected hole, lift this run through
its actual expansion inclusion and append the one contextual compute step.
The endpoints are $C[L]$ and $C[K]$, with the search token at the
corresponding host source position.

\lean{Copying/OperationalInputSafety}{177}{182}{search_run}

The theorem \texttt{consume} assembles the original residual specimen $D$.
Its context is $C'$, its families are the two remaining families, and its
token is the host image of the original entered source. Focus freedom follows
from plugging the focus-free right component into the original context.
Hole freedom on the left follows from preservation by the actual rewrite;
hole freedom on the right follows from the unchanged original right graph.
Transport these properties through the two expansion-peeling isomorphisms.
Term typing follows from the unchanged host-vertex label.

The two pointed isomorphisms then give equalities of representation classes.
Projecting the finite run proves
\[
 [S_L]\steps{c+1}[D_L],\qquad [D_R]=[S_R],
 \qquad c\in\{5,3\}.
\]
The token of $D$ is a search token and is not entering. Other original
shallow or deep paired components remain present with their original
certificates. The theorem is uniform in the chosen context class; its one
class-specific premise is that $C'$ belongs to that class. The final
input-safety theorems below discharge this premise.

\lean{Copying/OperationalInputSafety}{187}{283}{OperationalSafety.consume}

\subsection{Discharging the context and quantitative obligations}
For an arbitrary entering original specimen, its hole relation is recovered
from the selected edge certificate and global hole freedom. The input port
is recovered from actual source-list membership. The unique-input result
identifies it with the input used by the exact entry prefix.
Apply the preceding residual construction and choose the settled-specimen
case of input safety, with $c$ additional left steps after the first and
zero right steps. The resulting quantitative requirement is exactly
$Q(c+1,0)$. The equivalence-up-to parameter $P$ is arbitrary because this
case uses actual residual endpoints.

\lean{Copying/OperationalInputSafety}{287}{315}{inputSafe_of_plugging}

For the unrestricted context class, membership of $C'$ is immediate.
For the binding-free class, global hole freedom of the original right state
implies hole freedom of $H$; together with its focus freedom this makes $H$
binding-free. The entered hole has a source, so the already proved
sourced-hole plugging theorem applies. This does not assume the invalid
unrestricted closure statement for zero-input holes. Both final quantitative
theorems retain only the explicit arithmetic comparison $Q(c+1,0)$.

\lean{Copying/OperationalInputSafety}{321}{326}{template_inputSafe_cost}
\lean{Copying/OperationalInputSafety}{331}{339}{template_bindingFree_inputSafe_cost}

\termdef{operational-input-safety}
Specialising $Q$ to the universal cost relation proves input safety of both
operational families in each context class, with no remaining semantic
premise. For the binding-free class, the improvement relation also applies:
the right prefix has cost zero, which is no greater than the left cost.
This retains a quantitative comparison before costs are forgotten for
termination equivalence. Exact-cost input safety is not claimed here:
six or four transitions are not equal to zero.

\lean{Copying/OperationalInputSafety}{343}{346}{template_inputSafe}
\lean{Copying/OperationalInputSafety}{349}{352}{template_bindingFree_inputSafe}
\lean{Copying/OperationalInputSafety}{356}{359}{template_bindingFree_inputSafe_improvement}

\subsection{Original-specimen regression examples}
The isolated-specimen constructor places an actual closed operational pair
in the two fillings of one sourced term-hole corolla. Its output interface
is empty, and the context is binding-free. Focus freedom comes from the
pair; the supplied endpoint hole-freedom proofs make the complete fillings
hole-free. The token starts at the original term input. Applying
\texttt{consume} to this specimen and discharging sourced-hole plugging
gives a genuine settled residual with the exact prefix and unchanged right
representation class.

\lean{Copying/OperationalInputSafetyExamples}{13}{30}{Examples.isolated}
\lean{Copying/OperationalInputSafetyExamples}{34}{49}{isolated_consumed}

For reference creation, use the unchanged translated natural-number
constant as the stable argument. The operational pair is the literal
allocation rule with identity endpoint isomorphisms. Hole freedom of its
left side follows from serial assembly of the reference operation and
the translated constant; hole freedom of its right side follows from the
allocation construction. These proofs build the original specimen, and
the general consumption theorem supplies exactly four transitions to its
settled residual.

\lean{Copying/OperationalInputSafetyExamples}{51}{54}{referencePair}
\lean{Copying/OperationalInputSafetyExamples}{56}{59}{referenceLeft_closed}
\lean{Copying/OperationalInputSafetyExamples}{61}{63}{referenceSpecimen}
\lean{Copying/OperationalInputSafetyExamples}{67}{74}{reference_consumed}

For micro-beta, use the original translated variable body, including its
single-use distributor, with the same constant translation raised to the
required allowance. Stability survives raising. The actual application
rule gives the pair, and its proved safety laws give hole freedom on both
sides. The corresponding original binding-free specimen consumes its
entered hole in exactly six transitions.

\lean{Copying/OperationalInputSafetyExamples}{78}{78}{betaArgument}
\lean{Copying/OperationalInputSafetyExamples}{80}{81}{betaArgument_stable}
\lean{Copying/OperationalInputSafetyExamples}{83}{87}{betaPair}
\lean{Copying/OperationalInputSafetyExamples}{89}{94}{beta_sides_closed}
\lean{Copying/OperationalInputSafetyExamples}{96}{97}{betaSpecimen}
\lean{Copying/OperationalInputSafetyExamples}{100}{106}{beta_consumed}

Settlement is relative to the residual specimen context, not a claim that
execution has finished. In this beta example, the very next original
search transition requests copying from the retained variable-body
distributor. The checked transition confirms that consumption has not
silently performed the later copying or substituted the argument eagerly.

\lean{Copying/OperationalInputSafetyExamples}{110}{113}{beta_next_requests_copy}
